

\section{Conclusion}
\label{sec:conclusion}

We presented a \textbf{measurement framework for epidemiological model completeness} as a
systematic instrument for quantifying the gap between what a paper specifies and what can be
reconstructed from it as a formal compartment model.
LLMs served as simulated extractors for models described in epidemiological papers.
Phase~1 builds structured reference models
and paper promise inventories. Phase~2 employs LLM-assisted entity extraction from PDFs
across three commercial backends, evaluated with a domain-aware fuzzy string matching
protocol. Phase~3 closes extraction gaps via iterative retrieval-augmented completion,
structural validation, and rule-based repair, then assigns a five-component completeness
score that jointly measures gap reduction, reference agreement, fill traceability,
parameter accuracy, and structural integrity.

Applied to 10 peer-reviewed modeling papers across 10 infectious diseases, the pipeline
achieves a mean completeness score of \textbf{$74.0\%$}. Phase~3 raises mean structural recall
from $0.822$ (Phase~2 draft) to $0.965$ (winner mode), with all three Phase~3 modes improving
recall over the Phase~2 baseline. All 10 diseases show detectable structural errors in
their LLM-extracted models before repair; $44\%$ are resolved automatically.
RAG-only achieves the highest recall in 9 of 10 diseases and is the safest mode because
its fills are traceable to specific paper passages. LLM-only adds genuine value as a
coverage tool, but its fills are unverifiable and lower the traceability component of
the completeness score. The distinction between RAG-sourced and LLM-inferred fills is
itself a measurement outcome: a gap closed by LLM inference but not by RAG is direct
evidence of a specification failure in the paper.

These results quantify a consistent pattern: LLMs reliably extract compartment structure
but struggle with flow paths and parameters embedded in ODE derivations or referenced
only by symbolic expressions. The specification-validation gap is measurable and
non-trivial: even with the enhanced extraction pipeline using GPT-5.4-mini, Phase~2 drafts
miss on average $28\%$ of flows and $12\%$ of parameters relative to the best-case
Phase~2 draft. Phase~3 reduces those gaps substantially, achieving perfect recall for
all three entity types in 7 of 10 diseases.

%\paragraph{Future work.}
%Several concrete directions emerge directly from the failure cases described in this paper.

%\noindent\textbf{Expand the dataset.}
%The most immediate improvement is to include more papers per disease (currently one each).
%A single paper per disease means that a paper-specific writing quirk---parameters in
%figure captions (flu), composite ODE coefficients (TB), Bayesian-only values
%(Ebola)---dominates the per-disease score. Adding two or three papers per disease would
%separate pipeline limitations from individual paper choices, and would enable
%within-disease variance analysis and more robust mean completeness estimates.

%\noindent\textbf{Add figure and table extraction.}
%Several of the worst failures (flu, Ebola) stem from numeric values that exist only
%inside figures, supplementary tables, or figure captions. Integrating OCR-based figure
%parsing and structured table extraction (beyond Camelot's current scope) would close
%this entire class of gaps without any change to the downstream phases.

%\noindent\textbf{Improve RAG for multi-part definitions.}
%TB and Malaria show that parameters defined incrementally across multiple sections are
%not recoverable by single-passage retrieval. A multi-hop RAG strategy---where a
%retrieved passage can trigger further retrieval of referenced symbols---would let the
%system assemble composite definitions that no individual passage contains in full.

%\noindent\textbf{Enhance the evaluation protocol.}
%The current fuzzy string matching works well for standard compartment and flow names
%but still struggles with highly domain-specific notation (e.g., Bayesian posterior
%symbols, host--vector composite terms). Extending the synonym library and adding a
%disambiguation step for homonyms (``Susceptible adults'' vs.\ ``Susceptible children'')
%would improve recall precision and reduce false matches that inflate precision scores.

%\noindent\textbf{Extend the structural validator.}
%Adding dimension-aware validation (e.g., detecting a rate parameter with units
%$\text{person}^{-1}$ assigned to a flow expecting $\text{person}^{-1}\text{day}^{-1}$)
%and a formal parameter-to-flow assignment step (to resolve orphaned parameters in
%multi-group models like HIV and Malaria) would significantly improve the structural
%integrity component, currently the weakest in the completeness score.

%\noindent\textbf{Test generalizability.}
%Applying the framework to stochastic, agent-based, and network models beyond ODE
%compartmental structures would test whether the three-phase approach and the completeness
%score translate to a broader class of mathematical models in infectious disease
%epidemiology. This also includes models where multiple interacting patches or age-strata
%appear as structured population layers rather than separate compartments.

%\noindent\textbf{Close the loop with simulation.}
%Finally, studying whether completeness scores correlate with downstream simulation
%accuracy---i.e., whether a higher-scoring extracted model reproduces the paper's reported
%trajectories more faithfully---would validate the score as a predictor of scientific
%reproducibility and not just a structural inventory metric.
